% Evaluation fixture for manuscript-reporting-check: a fictional tool paper prepared for double-blind
% review. The tool, the repositories, the Zenodo record and the timings are invented for testing.
\documentclass{article}
\usepackage{url}
\title{varpanel: assembling gene-panel reports from annotated VCFs}
\author{Anonymous}
\begin{document}
\maketitle

\begin{abstract}
Clinical research groups often re-annotate the same gene panels for every cohort. varpanel turns an
annotated VCF and a panel definition into a per-sample report in one command. On ten public exomes it
produced reports in a median of 38 seconds per sample on a laptop. The software is open source and
available at \url{https://anonymous.4open.science/r/varpanel-3F1C} (anonymised for review).
\end{abstract}

\section{Introduction}
Gene-panel reports are rebuilt by hand in many laboratories. varpanel reads a panel definition in YAML
and an annotated VCF, and writes an HTML and a TSV report per sample.

\section{Implementation}
varpanel is written in Python. Panel definitions are validated against a JSON schema before a run, and
the annotation fields to report are declared in the panel file.

\section{Evaluation}
We ran varpanel on ten public exomes deposited at Zenodo (record 7654321) with a 120-gene cardiomyopathy
panel. Table~\ref{tab:tools} compares it with two other tools.

\begin{table}[h]
\caption{Tools that build panel reports.}
\label{tab:tools}
\begin{tabular}{lll}
Tool & Access & Licence \\
varpanel & local command line & MIT \\
PanelServe & web server (by request) & proprietary \\
\end{tabular}
\end{table}

\section{Code and data availability}
\begin{itemize}
    \item \textbf{varpanel}: \url{https://anonymous.4open.science/r/varpanel-3F1C} (MIT).
    \item \textbf{Panel definitions and benchmark scripts}:
    \url{https://anonymous.4open.science/r/varpanel-panels-9A2E}.
    \item The ten exomes are public at Zenodo (record 7654321).
\end{itemize}

\appendix
\section{Panel schema}
Panel files declare genes, transcripts and report fields. Field descriptions are loaded into the
report template on demand, so a panel can add fields without changing the code.

\end{document}
